% Part III-B, introduction and related work (included by main_en.tex).

\section{Introduction}\label{sec:intro}

RLVR assumes that every sampled solution can be checked. When checking is expensive, for example with human grading,
long test suites or slow simulators, a training run may afford labels for only some of its samples. A judge can fill in
the rest. A judge that was right on yesterday's samples may be wrong on today's, though, because the policy it judges
keeps changing.

This paper studies that setting with GRPO on MATH, using Part II's setup, where the verifier labels 2 of the 8
completions in each group. We compare four ways to use the other 6: drop them, reward them with \kev{}'s frozen
judgment, reward them with \kev{} after its head is refitted online on the labels collected so far, or reward them with
a head of the same shape trained from scratch on the same labels. Because the verifier is in fact available, we also
measure how good each fill-in reward was on the completions it scored, step by step.

\section{Related work}\label{sec:related}

Semi-supervised reward learning and active labelling reduce the number of human labels in RLHF~\cite{ouyang2022training}.
Reward models trained once and then frozen are exploited as the policy drifts~\cite{gao2023scaling}. Re-fitting the
reward model to fresh labels during training is the usual remedy in RLHF, but it costs labels. Recalibrating a
classifier's output layer on new data is a standard fix for distribution shift~\cite{guo2017calibration}. Here it is
applied online to a judge's pointer head. GRPO and our update follow Part II~\cite{shao2024deepseekmath,liu2026tricks}.
